\section{Introduction}
Public leaderboards for LLM agents, from SWE-bench \citep{jimenez2024swebench} to tau-bench \citep{yao2024taubench} and the Holistic Agent Leaderboard \citep{kapoor2026hal}, rank systems by a single success rate obtained from one run per task. Claims such as ``new state of the art'' are made on differences of one or two points. Three sources of randomness are rarely separated: \emph{which tasks} were sampled, \emph{which run} of a stochastic agent was recorded, and \emph{how} the agent and the task interact. Work on error bars for language-model evaluations \citep{miller2024error,bowyer2025clt}, on variance in benchmarks \citep{madaan2024variance,bouthillier2021variance} and on reliability metrics for agents \citep{yao2024taubench,kapoor2024matter} has established the principle. What is missing is an agent-specific toolkit that (i) treats the task as the sampling unit while exploiting repeated runs, (ii) decomposes variance across agents, tasks, their interaction and run noise, (iii) tells a practitioner how to split a budget between more tasks and more repeats, and (iv) is validated against ground truth.

\paragraph{Hypothesis.} \emph{Many differences between agents on published leaderboards are not statistically distinguishable, and metrics such as pass@$k$ hide unreliability that pass$^k$ reveals.} We treat this as a hypothesis to test, not an assumption, and report where the data support it and where they do not.

\paragraph{Contributions.}
\begin{enumerate}
\item \textbf{\textsc{AgentAudit}}, a typed, tested Python library with a CLI and a browser playground that runs the same code (Pyodide). Statistics are implemented from scratch and cross-checked against SciPy and statsmodels (Section~\ref{sec:method}).
\item \textbf{\textsc{AgentAudit-Bench}}: six offline deterministic environments, 360 tasks with horizon and distractor knobs, verification by final state, seeded and logged tool failures, and record/replay caching (Section~\ref{sec:bench}).
\item \textbf{An empirical study} with three legs: simulations with agents of known skill (coverage, error rates, power, allocation, sequential stopping, ranking calibration); a re-analysis of public SWE-bench Verified and tau-bench results; and \REALINTRO (Section~\ref{sec:results}).
\item \textbf{Practical rules} derived from the variance decomposition: a design effect $1+(R-1)\rho$ for pooled intervals, a closed form for the cost-optimal number of runs per task, and a quantitative comparison of confidence sequences with fixed-sample tests and naive peeking.
\end{enumerate}

\paragraph{Findings in brief.} (1) Pooling runs is not merely loose: ten runs per task push the coverage of a nominal 95\% pooled Wilson interval to \covWilsonTenRuns{}, while task-level bootstrap and hierarchical beta-binomial intervals stay within \covBootRange. (2) Pooled run-level tests have type-I error of \typeIPoolMc{} (nominal 5\%) when agent-by-task interaction is large; task-level sign-flip tests stay in \typeIPermMin--\typeIPermMax. (3) At a fixed budget of \designBudget{} attempts, power falls from \designPowerOne{} with one run per task to \designPowerTen{} with ten. (4) On SWE-bench Verified, \sweIndistHolm{} of pairs are indistinguishable and no adjacent pair is separated. (5) A confidence sequence controls false positives where naive peeking does not (\seqNullSeq{} versus \seqNullPeek{} under the null) but costs power against a fixed-sample test at the same budget. \REALFINDINGS
